\section{NP-hardness for classical mass action}
\label{sec:hardness-independent}

\begin{theorem}[Weak NP-hardness]
\label{thm:hardness-independent}
Weak all-mobile stationary-Turing admissibility of rationally encoded finite
classical mass-action networks is NP-hard under polynomial-time many-one
reductions from \textsc{Partition}. For an input of length $\ell$, the
construction chooses $m=k^2>\ell$ and has $3m+1$ species and $8m+2$ indexed
reactions. Every stoichiometric column is $\pm e_i$, the stoichiometric rank
is full, and all rates and equilibrium coordinates in a YES witness are
strictly positive.
\end{theorem}

\begin{proof}
Form the lifted matrix family $L(q)$ from the open-cube instance. It has
$m+1$ fixed rows and $2m$ variable rows. Apply
Theorem~\ref{thm:row-realization-independent}; the complete positive
steady-flux image is
\[
 A(v)=R L(q),\qquad R\succ0\text{ diagonal},\quad q\in(-1,1)^{2m}.
\]
At an arbitrary positive equilibrium, every Jacobian is
\[
 J=R L(q)H,
 \qquad H\succ0\text{ diagonal}.
\]
Conjugation by $R$ gives
\[
 R^{-1}JR=L(q)HR=L(q)D,
 \qquad D=HR\succ0.
\]
The scaling-elimination theorem therefore says that some positive-equilibrium
Jacobian is Hurwitz exactly when the partition instance is YES; no theorem
about unrestricted diagonal stabilizability is invoked.

The coordinate inherited from the bottom-right entry of the open-cube matrix
is a fixed row and column of the lift, with $L_{jj}=k\beta>0$. Hence every
realized Jacobian has
\[
 J_{jj}=\rho_jk\beta h_j>0,
\]
so its signed singleton minor is $-J_{jj}<0$. Whenever the homogeneous
Jacobian is Hurwitz, the fixed-J diagonal-damping theorem supplies a strictly
positive diffusion matrix producing a positive real nonzero-mode eigenvalue.
Conversely, admissibility includes homogeneous stability. Thus the reaction
network is stationary-Turing-admissible exactly when the partition instance
is YES.

For a YES instance choose the rational interior witness of the open-cube
proof, $R=H=I$, $x^*=\mathbf1$, and rates equal to the strictly positive
row-realization fluxes. There are
\[
 2(m+1)+3(2m)=8m+2
\]
reactions. Denominator clearing and the base-complex construction have
polynomial bit complexity. This is a polynomial-time many-one reduction.
\end{proof}

The theorem is ordinary (weak) NP-hardness because the source problem is
\textsc{Partition}. The source and target complexes produced by denominator
clearing can have molecularity growing with the input; no bounded-molecularity
hardness statement follows from this construction.
